\chapter{视觉模型}
\begin{quote}\small
\textit{本章摘要。}本章从图像的局部空间结构出发，说明卷积为何形成层级表示、残差为何改善深层优化，再逐步扩展到检测、分割、Vision Transformer 和自监督学习。最后将评价从单张图像分数推进到跨相机、跨批次、拒识和数据追溯，强调视觉模型必须服务于可验收的工业决策。
\end{quote}
风险预警告诉维护人员何时需要关注设备，却不能指出缺陷位于哪里。本章转向巡检图像和表面质检：先判断是否存在缺陷，再定位区域，最后估计边界与尺寸。这些任务使用图像级、框级或像素级标签，与未来故障标签不同。拆检后才能获得的照片可以支持离线诊断，却不能直接作为拆检前预警的输入。
上一章中，时间顺序决定同一组读数代表稳定运行还是异常演化；本章中，空间关系决定相同的颜色和纹理代表正常表面还是裂纹。把图像展平成无结构向量虽然在数学上可行，却放弃了邻近像素之间的规律和不同位置共享的成像机制。卷积把这种共享写进参数化，注意力则让远距离区域按内容交互，两者是在不同数据和资源条件下组织空间证据的方式。

识别“这里有裂纹”和指出裂纹的准确位置，对图像细节的要求不同。下采样能减少计算、扩大观察范围，但一条很细的裂纹也可能随之消失。即使网络能比较图像中相距很远的部件，也无法恢复输入时已经丢掉的细节。因此先确定最小要检出的缺陷有多大、需要输出类别还是位置，再选择分辨率和网络。下面从卷积和感受野推到检测框、像素损失与校准，逐步说明这些选择怎样影响结果。
\section{视觉表示：从卷积骨干到视觉 Transformer}

% auxiliary-resource-guide
本章末的“辅助学习材料”列出与核心方法对应的讲解和实践入口，可在阅读相关推导后，按所列小节对照学习。

Residual Network (ResNet) 通过残差连接改变深层网络的优化路径，ViT 将图像块编码为 token 后使用 Transformer\cite{he2016resnet,dosovitskiy2021vit}。下文从这两种机制比较局部归纳偏置、全局交互和预训练需求。
\begin{figure}[htbp]
\centering
\begin{tikzpicture}[node distance=0.65cm and 0.65cm]
 \node[idi input, minimum width=1.7cm] (image) {工业图像\\$H\times W$};
 \node[idi process, right=of image, minimum width=2.0cm] (backbone) {视觉骨干\\CNN / ViT};
 \node[idi state, right=of backbone, minimum width=2.0cm] (features) {多尺度特征\\$F_1,\ldots,F_s$};
 \node[idi output, right=of features, minimum width=1.8cm] (class) {分类\\图像级};
 \node[idi output, below=0.75cm of class, minimum width=1.8cm] (detect) {检测\\框级};
 \node[idi output, above=0.75cm of class, minimum width=1.8cm] (segment) {分割\\像素级};
 \draw[idi arrow] (image)--(backbone); \draw[idi arrow] (backbone)--(features);
 \draw[idi arrow] (features)--(class); \draw[idi arrow] (features)--(detect); \draw[idi arrow] (features)--(segment);
\end{tikzpicture}
\caption{视觉系统先由卷积或 Transformer 骨干提取表示，再按标签粒度输出分类、检测或分割结果。}
\label{fig:vision-output-granularity}
\end{figure}
图\ref{fig:vision-output-granularity}将共享骨干与三种输出粒度分开：分类使用图像级标签，检测和分割还要求保留位置与边界。图像是具有二维邻域结构的信号。卷积层以局部窗口提取平移共享的特征：
\begin{equation}
 z_{i,j,c}=\sum_{u,v,c'}W_{u,v,c',c}x_{i+u,j+v,c'}+b_c.
\end{equation}
共享权重降低参数量，也把局部性作为归纳偏置。堆叠卷积、非线性和下采样后，感受野逐步扩大，表示从边缘、纹理发展到部件和对象。

\paragraph{卷积的索引、参数与等变性。}
设输入$X\in\R^{H\times W\times C_{\rm in}}$，方形核大小$k$，输出通道数$C_{\rm out}$，步幅$s$、填充$p$、空洞率$d_a$。采用从零开始的空间索引，有
\begin{equation}
 Z_{i,j,c}=b_c+\sum_{u,v=0}^{k-1}\sum_{a=1}^{C_{\rm in}}
 K_{u,v,a,c}X_{si-p+d_au,\,sj-p+d_av,\,a},
\end{equation}
越界输入按指定规则填充。这个深度学习中惯称的卷积实际采用互相关索引；若使用数学卷积，只需翻转核。为使最右端采样不超过填充后的输入，需满足$s i+d_a(k-1)\leq H+2p-1$，因此
\begin{equation}
 H_{\rm out}=\left\lfloor\frac{H+2p-d_a(k-1)-1}{s}\right\rfloor+1,
\end{equation}
宽度同理。参数量为$k^2C_{\rm in}C_{\rm out}+C_{\rm out}$，与图像面积无关；不共享的局部连接则须为每个输出位置另存一组权重。

在无限网格或忽略边界、$s=1$时，令$(T_\delta X)_{i,j}=X_{i-\delta_1,j-\delta_2}$，直接代入求和式得$K*(T_\delta X)=T_\delta(K*X)$。这就是平移等变性，而不是平移后输出保持不变。步幅$s>1$时通常仅保留对$s$整数倍平移的对应关系，边界填充还会引入额外差异。图像分类可以通过汇聚追求不变性，定位任务却应保留位置随输入一起移动的关系。

深度可分离卷积先对每个输入通道施加一个空间核，再用逐点矩阵混合通道。忽略偏置，每个输出位置的乘加数由$k^2C_{\rm in}C_{\rm out}$降为$k^2C_{\rm in}+C_{\rm in}C_{\rm out}$，比值为$1/C_{\rm out}+1/k^2$。节省来自结构约束：对固定输入通道，通向各输出通道的空间模板被限制为同一个模板的不同倍数，因此它不等价于任意标准卷积，且实际延迟还受内存访问影响。

卷积骨干的演化可以理解为三个连续问题：如何让网络更深而仍然可优化，如何在有限计算量下提高特征复用，如何让不同尺度的目标都得到足够的空间分辨率。ResNet通过残差连接$h_{l+1}=h_l+F_l(h_l)$改善深层网络的优化；DenseNet进一步把前面各层的特征拼接到当前层，使梯度和低层纹理可以直接复用，但显存访问成本更高。MobileNet 使用深度可分离卷积，将标准卷积拆成逐通道空间卷积和$1\times1$逐点卷积，计算量近似从$k^2C_{in}C_{out}$降为$k^2C_{in}+C_{in}C_{out}$，适合边缘质检；EfficientNet 则联合缩放深度、宽度和输入分辨率，强调在给定算力下平衡三者，而不是单独堆叠层数。

对工业图像而言，骨干网络并不只是分类器的前端。ResNet-18/34 通常适合小数据和低延迟基线，ResNet-50 适合迁移学习的性能--成本折中，ConvNeXt 把大卷积核、归一化和训练策略现代化，在保留卷积局部归纳偏置的同时接近 Transformer 的表示能力。选择骨干时应同时记录参数量、特征图步幅、显存、冻结层数和输入分辨率；同一个“准确率更高”的模型，如果把最小可检缺陷下采样掉，就不适合作为生产模型。

ViT 把图像划分为$P\times P$的 patch，展平后通过线性投影得到 token：$z_i=x_iW_E+e_i^{pos}$，再用标准 Transformer 建模全局关系。其优势是任意两个区域可以在少量层数内直接交互，代价是 patch 数量为$N=HW/P^2$时注意力开销近似为$O(N^2d)$，并且从零训练往往需要更多数据。Swin Transformer 通过窗口注意力和交错移位把复杂度降为局部窗口规模的平方，并通过层级下采样适配检测和分割；当缺陷同时依赖局部纹理和较大范围工件结构时，Swin 或 ConvNeXt 往往比原始 ViT 更自然。

卷积与 Transformer 不是简单的替代关系。卷积提供局部性、平移等变性和较强的数据效率，ViT 提供长距离交互和可扩展的全局建模。工业迁移实验应固定预训练数据、输入分辨率和微调预算，至少比较从零训练、冻结骨干线性探测、全量微调和参数高效微调四种设置，避免把预训练规模或训练预算误认为模型结构的收益。

\begin{table}[htbp]
\centering
\caption{视觉模型族谱与工业选型依据}
\label{tab:vision-model-families}
\scriptsize
\setlength{\tabcolsep}{3pt}
\resizebox{\linewidth}{!}{%
\begin{tabular}{p{2.4cm}p{3.7cm}p{3.5cm}p{4.0cm}}
\hline
模型族 & 核心机制 & 主要优势 & 典型工业选择依据 \\
\hline
ResNet & 残差块与层级卷积 & 优化稳定、预训练模型与工具较完善 & 小中型数据集、稳定基线、低风险迁移 \\
DenseNet & 跨层特征拼接 & 特征复用充分、参数相对紧凑 & 纹理细节重要但显存带宽充足 \\
MobileNet/EfficientNet & 深度可分离卷积或复合缩放 & 参数和延迟较低 & 边缘设备、相机侧实时推理 \\
ConvNeXt & 现代化大卷积核骨干 & 局部先验与表达能力平衡 & 需要卷积归纳偏置的高性能分类或检测 \\
ViT/Swin & patch 注意力或窗口注意力 & 全局关系与大规模预训练迁移 & 长距离结构、充足预训练或多尺度输出 \\
\hline
\end{tabular}
}
\end{table}

表\ref{tab:vision-model-families}把结构机制对应到工业资源条件，比较时应区分残差优化、特征复用、边缘计算和长距离交互各自解决的问题。
\paragraph{表示学习的评价。}
分类常用top-1准确率，检测常用平均精度，分割常用交并比。工业视觉还应报告跨设备、跨光照、跨批次和少见缺陷的结果。视觉模型的显著性图可以辅助分析，但可视化不等于因果解释，仍需结合遮挡实验、反事实扰动或专家标注验证。

\paragraph{卷积的空间归纳偏置。}
卷积核是局部加权模板。平滑核抑制高频噪声，边缘核响应局部变化，多个卷积核共同形成一组可学习滤波器。步幅大于一会降低空间分辨率，空洞卷积在不显著增加参数的情况下扩大感受野。池化提供一定平移鲁棒性，却也可能丢失小缺陷位置，因此检测任务常使用特征金字塔恢复多尺度信息。

\paragraph{残差优化与特征复用。}
普通深层网络学习$H(x)$，残差网络学习$F(x)=H(x)-x$。若恒等映射已经足够好，优化器只需令$F$接近零。对残差块$h_{l+1}=h_l+F_l(h_l)$，梯度包含直接项$I$，因此比连续相乘的纯变换更容易沿深层传播。残差连接也成为Transformer和扩散模型的基础组件。

\paragraph{残差梯度不再只有单一路径。}
设各层$h_l\in\R^d$，且残差分支$F_l:\R^d\to\R^d$可微。令$J_l=\partial F_l/\partial h_l^{\mathsf T}$，则
\begin{align}
 \frac{\partial h_{l+1}}{\partial h_l^{\mathsf T}}&=I+J_l,\\
 \nabla_{h_l}\mathcal L&=(I+J_l)^{\mathsf T}\nabla_{h_{l+1}}\mathcal L.
\end{align}
连续两个块的前向Jacobian为$(I+J_{l+1})(I+J_l)=I+J_l+J_{l+1}+J_{l+1}J_l$，包含不经过任何残差分支的恒等项、经过一个分支的项和经过两个分支的项。相比普通网络只有$J_{l+1}J_l$，残差结构允许梯度走较短路径。若$\|J_l\|_2\leq\epsilon<1$，则对任意向量$v$有
\begin{equation}
 (1-\epsilon)\|v\|_2\leq\|(I+J_l)v\|_2\leq(1+\epsilon)\|v\|_2.
\end{equation}
小残差分支使单块接近恒等，但多块的上下界仍相乘，所以残差不是绝不消失或爆炸的定理。改变通道或空间尺寸时，要把捷径替换为投影$S_lh_l$，导数变为$S_l+J_l$，此时也不再有字面上的恒等路径。

\paragraph{图像块与窗口注意力的成本。}
假设$H,W$均可被$P$整除，令第$i$个展平图像块为列向量$x_i\in\R^{P^2C}$，投影矩阵$E\in\R^{d\times P^2C}$，则$z_i=Ex_i+e_i^{\rm pos}\in\R^d$；前文$x_iW_E$采用的是其转置后的行向量记法。共有$N=HW/P^2$个块，$QK^{\mathsf T}$需要$O(N^2d)$计算。若每个窗口包含$M\times M$个token，共有$N/M^2$个窗口，局部注意力开销为
\begin{equation}
 O\left(\frac{N}{M^2}(M^2)^2d\right)=O(NM^2d).
\end{equation}
固定$M$时对$N$线性增长，但单层不能直接连接不同窗口，移位窗口用后续层交换边界信息。另一方面，线性patch投影若$d<P^2C$，其秩最多为$d$，必有非零输入变化被映射为零；这说明不可假定所有细节都保留在token中，却不意味着小于patch的缺陷必然不可见，是否丢失还取决于投影和训练。

\section{空间输出：检测、分割与多尺度特征}
卷积神经网络（Convolutional Neural Network，CNN）可以用于不同粒度的视觉预测。U-Net 的同尺度跳跃连接用于恢复像素级定位，Feature Pyramid Network (FPN) 通过特征金字塔组织多尺度语义，Mask Region-based Convolutional Neural Network (Mask R-CNN) 则在实例检测中增加掩码分支\cite{ronneberger2015unet,lin2017fpn,he2017mask}。这些方法解决的是不同粒度的空间输出，不能仅按网络深度排序。
从图像分类转向定位时，高层语义与空间细节之间的矛盾变得直接。高分辨率特征保留小缺陷边界，低分辨率特征提供较大范围的上下文。特征金字塔通过自顶向下路径与横向连接融合不同尺度，因此应先根据缺陷在输入图像中的像素尺寸选择输出层级，再选择检测头。对于 ViT，patch 过大同样可能丢失细缺陷，卷积 stem、重叠 patch 和层级表示可以提供不同折中。小、中、大目标的分层召回有助于区分表示分辨率不足与后续匹配错误。

检测和分割的差异不只是输出张量的形状，而是模型必须建立从候选区域到空间决策的完整机制。两阶段检测器 Faster R-CNN 先由区域建议网络产生候选框，再由 感兴趣区域（Region of Interest，RoI）的池化或 RoIAlign 提取每个候选区域的特征并分别分类、回归；它通常具有较高定位质量，但候选区域和多阶段计算会增加延迟。单阶段 You Only Look Once (YOLO)、RetinaNet 和 Fully Convolutional One-Stage Object Detection (FCOS) 直接在密集特征图上预测类别与位置，速度更快，却必须处理大量背景位置和正负样本分配。RetinaNet 的 focal loss
\begin{equation}
 \mathcal{L}_{\mathrm{focal}}=-\alpha_t(1-p_t)^\gamma\log p_t
\end{equation}
其中$t$表示该样本的真实类别，$p_t$是模型对真实类别的预测概率，$\alpha_t$是该类别的预设权重，$\gamma\geq0$控制对易分类样本的降权。二分类时，$p_t$按真实标签在$p$与$1-p$之间选取。
通过降低易分类背景的权重缓解类别不平衡；FCOS 不依赖预设 anchor，直接回归位置到边界的距离，减少 anchor 尺寸和比例的超参数。工业小缺陷任务应比较 anchor-based、anchor-free 和不同输入分辨率，而不能只更换 backbone。

\paragraph{检测框、匹配与定位损失。}
检测器为每个候选框输出类别概率和边界框参数。两个框$A,B$的交并比为
\begin{equation}
 \operatorname{IoU}(A,B)=\frac{|A\cap B|}{|A\cup B|}.
\end{equation}
训练通常由分类损失和定位损失组成。推理时非极大值抑制保留高分框并删除与其重叠过大的框。小缺陷检测需要特别关注输入分辨率、正负样本比例和标注框的一致性。

\paragraph{从几何框到可训练的回归目标。}
设参考框中心及宽高为$(a_x,a_y,a_w,a_h)$，真实框为$(g_x,g_y,g_w,g_h)$，宽高均严格为正且使用同一像素坐标系。直接回归像素位移会使大框产生更大尺度的梯度，因此定义无量纲目标
\begin{equation}
 t_x^*=\frac{g_x-a_x}{a_w},\quad t_y^*=\frac{g_y-a_y}{a_h},\quad
 t_w^*=\log\frac{g_w}{a_w},\quad t_h^*=\log\frac{g_h}{a_h}.
\end{equation}
网络输出$t\in\R^4$后，以$b_x=a_x+a_wt_x,b_y=a_y+a_ht_y,b_w=a_we^{t_w},b_h=a_he^{t_h}$解码。中心相对化降低尺度依赖，对数宽高保证解码尺寸为正；但$\partial b_w/\partial t_w=b_w$，极端输出仍可能带来不稳定。

对匹配到真实框的正样本，常用残差$e=t-t^*$的分段损失
\begin{equation}
 \rho_\beta(e)=\begin{cases}e^2/(2\beta),&|e|<\beta,\\|e|-\beta/2,&|e|\geq\beta,\end{cases}
 \qquad \beta>0.
\end{equation}
其导数在小误差处为$e/\beta$，大误差处为$\operatorname{sign}(e)$，兼顾局部平滑与异常残差的有界影响。完整损失可写为
\begin{equation}
 \mathcal L=\frac1{N_{\rm cls}}\sum_i\ell_{\rm cls}(p_i,y_i)
 +\frac{\lambda}{\max(1,N_+)}\sum_{i:y_i>0}\sum_{j\in\{x,y,w,h\}}\rho_\beta(t_{ij}-t_{ij}^*).
\end{equation}
背景样本没有真实框回归目标；分类与定位的归一化方式必须报告，否则相同$\lambda$可能代表不同梯度比例。

IoU直接比较几何重叠。设两个边长为$a$的正方形仅水平错位$\delta\in[0,a)$，则交集为$a(a-\delta)$，并集为$a(a+\delta)$，所以$\operatorname{IoU}=(a-\delta)/(a+\delta)$，在$\delta=0$右侧的导数为$-2/a$。这说明相同一像素偏移对小框更加苛刻。不相交框的普通IoU为零，在保持不相交的区域内不给出靠近方向，因此实际定位还会结合参数回归或加入外接框、中心距离等项。

\paragraph{焦点损失究竟改变了什么梯度。}
对正类样本设logit为$z$、$p=\sigma(z)$，则$\ell=-\alpha(1-p)^\gamma\log p$。不能只把交叉熵梯度乘以$(1-p)^\gamma$，因为这个权重本身依赖参数。乘积法则给出
\begin{equation}
 \frac{d\ell}{dz}=\alpha\gamma p(1-p)^\gamma\log p-\alpha(1-p)^{\gamma+1}.
\end{equation}
令$p=1-\varepsilon$并利用$\log(1-\varepsilon)\approx-\varepsilon$，得梯度约为$-\alpha(\gamma+1)\varepsilon^{\gamma+1}$；普通交叉熵则约为$-\alpha\varepsilon$。所以易样本梯度以更高阶衰减，困难样本占据更大的相对权重。但标注错误往往也是困难样本，焦点损失不具备自动甄别错误标签的能力，也不能保证输出概率仍然校准。

\paragraph{语义分割、实例分割与 U-Net。}
像素级交叉熵适合类别互斥场景，Dice损失
\begin{equation}
 \mathcal{L}_{Dice}=1-\frac{2\sum_i p_iy_i+\epsilon}{\sum_i p_i+\sum_i y_i+\epsilon}\label{eq:vision-soft-dice}
\end{equation}
对小目标更友好。U-Net 用逐级下采样的编码器提取语义上下文，再用逐级上采样的解码器恢复空间分辨率，并通过同尺度跳跃连接把编码器中的细粒度特征送入解码器。跳跃连接缓解了下采样造成的边界丢失，但也会把纹理噪声带入输出，因此需要结合标注质量、通道宽度和后处理验证。U-Net++通过重新设计跳跃路径减少编码器与解码器之间的语义鸿沟，DeepLab则以空洞空间金字塔池化扩大多尺度上下文；它们解决的是不同的空间信息瓶颈，不能仅凭网络名称比较优劣。实例分割还要进一步区分相邻目标，Mask R-CNN 在检测框分支之外增加 RoI 级掩码分支。边界质量、缺陷面积误差和漏检区域的连通性可能比平均 IoU 更接近工业验收标准。对关键缺陷，可以使用多个模型或随机增强估计预测波动。

\paragraph{像素似然、Dice与稀疏前景。}
设一幅图含$n$个像素，二值标签$y_i\in\{0,1\}$，模型输出$p_i=\sigma(z_i)$。在给定图像后用逐像素Bernoulli分布作为训练分解，有
\begin{align}
 p(y\mid X)&=\prod_{i=1}^n p_i^{y_i}(1-p_i)^{1-y_i},\\
 \mathcal L_{\rm BCE}&=-\frac1n\sum_i[y_i\log p_i+(1-y_i)\log(1-p_i)],\\
 \frac{\partial\mathcal L_{\rm BCE}}{\partial z_i}&=\frac{p_i-y_i}{n}.
\end{align}
这个似然分解不表示真实裂纹像素相互独立，而是输出头采用的近似。前景极少时，大量背景项会主导平均损失。Dice从集合重叠出发：硬掩码的$D=2|A\cap B|/(|A|+|B|)$，将预测指示函数换成概率便得到式\eqref{eq:vision-soft-dice}中的软Dice。记$N=2\sum_i p_iy_i+\epsilon$、$D_0=\sum_i p_i+\sum_i y_i+\epsilon$，则
\begin{equation}
 \frac{\partial\mathcal L_{\rm Dice}}{\partial p_j}
 =\frac{N-2y_jD_0}{D_0^2},\qquad
 \frac{\partial\mathcal L_{\rm Dice}}{\partial z_j}
 =\frac{N-2y_jD_0}{D_0^2}p_j(1-p_j).
\end{equation}
一个像素的梯度依赖全图预测总量，体现其区域级归一化，而不是各像素独立计分。对无前景图，$N=\epsilon$，当预测面积较大时惩罚梯度可能很小；可结合Binary Cross Entropy (BCE)并明确空掩码的评价约定。Dice也不是概率校准目标，因此不应把优化后的$p_i$未经验证就解释为可靠风险概率。

硬掩码下，令$J=\operatorname{IoU}$，由$|A|+|B|=|A\cup B|+|A\cap B|$得到$D=2J/(1+J)$。同一对掩码的两个指标单调对应，但先对图像平均再作非线性变换不等价；软Dice与阈值后二值IoU更不能直接替换。边界和拓扑验收正是补充这些面积指标遗漏的信息。

\paragraph{数据增强与标签语义。}
裁剪、翻转、颜色扰动和噪声增强应保持标签语义不变。若翻转会改变设备方向，或颜色变化会掩盖真实缺陷，增强反而制造错误先验。增强策略应由物理过程和相机成像条件约束，并通过消融实验验证。

\paragraph{工业检测案例。}
对表面缺陷数据，可先建立“人工特征加树模型”和“小型卷积网络”两个基线，再比较预训练视觉编码器。实验应按生产批次或相机划分，报告少数缺陷的召回率和每千张图的误报数。IoU 对边界移动的敏感性、正负样本不平衡的来源以及不改变缺陷标签的增强边界，都应作为验收规则的一部分明确记录。

\paragraph{感受野与有效感受野。}
理论感受野描述某个深层单元能够访问的输入区域。若连续两层卷积核大小为$k_1,k_2$、步幅为$s_1,s_2$，感受野可递推为
\begin{equation}
 r_2=r_1+(k_2-1)j_1,\qquad j_2=j_1s_2,
\end{equation}
其中$j$表示相邻深层单元在原图上的间隔。更一般地，令$r_0=j_0=1$，第$l$层核大小、步幅与空洞率为$k_l,s_l,d_l$。该层最左与最右采样点间隔为$d_l(k_l-1)$个上一层步距，每个采样点又覆盖$r_{l-1}$个输入位置，故
\begin{equation}
 r_l=r_{l-1}+d_l(k_l-1)j_{l-1},\qquad j_l=s_lj_{l-1}.
\end{equation}
例如三层$3\times3$卷积的步幅依次为$1,2,1$、空洞率均为一，则各层的感受野与步距为
\begin{align*}
 (r_1,j_1)&=(3,1),\\
 (r_2,j_2)&=(5,2),\\
 (r_3,j_3)&=(9,2).
\end{align*}
第三层相邻位置在原图相隔两像素，各自理论上覆盖九像素宽的区域。这里计算的是内点的覆盖范围，不保证边界处拥有同样多的真实像素，也不保证这个范围内每个像素都获得显著权重。理论感受野很大不代表所有像素贡献相同，实际梯度通常集中在更小的有效区域。小缺陷任务需要检查下采样是否在进入高层前已经抹去目标。

\paragraph{等变性、池化与空间细节。}
理想卷积具有平移等变性：输入平移后，输出特征也相应平移。但边界填充、步幅和池化会破坏严格等变性。池化提供一定不变性，却可能让两个相邻缺陷无法区分。工业视觉中应根据缺陷尺寸选择下采样比例，并在增强实验中单独测试平移、旋转、尺度和光照变化。

\paragraph{归一化层与批次统计。}
Batch Normalization 使用批内均值和方差，训练与推理阶段的统计量来源不同；小批量或设备分布高度异质时，估计会很噪。Layer Normalization 按样本内特征归一化，更适合序列和 Transformer。Group Normalization 在视觉小批量训练中常更稳定。归一化不是无害的数值操作，它改变了网络的参数化和对批次组成的依赖。

\paragraph{检测匹配与正样本分配。}
检测训练首先要决定预测框与真实框如何匹配。基于交并比（Intersection over Union，IoU）的正负样本分配简单，但小目标和稀疏目标可能没有足够正样本；一对多匹配会导致重复预测，需配合非极大值抑制或集合匹配。定位损失的尺度也影响分类与回归梯度的平衡，应报告损失权重和正样本定义。

\paragraph{分割边界与拓扑。}
两个掩码具有相同 IoU，不代表它们对工业决策同样有用。细长裂纹、孔洞和连通结构需要边界距离、连通性和骨架指标。Dice 对小目标友好，但可能忽略边界位置；边界损失可以补充几何信息。训练标签若边缘本身含糊，模型输出精细边界并不一定是更真实的结果。

\begin{figure}[htbp]
\centering
\begin{tikzpicture}[>=Latex, every node/.style={font=\small}]
 \draw[fill=blue!8, draw=blue!60] (0,0) rectangle (2.4,2.4);
 \foreach \x in {0.4,0.8,1.2,1.6,2.0} {\draw[gray!35] (\x,0)--(\x,2.4);}
 \foreach \y in {0.4,0.8,1.2,1.6,2.0} {\draw[gray!35] (0,\y)--(2.4,\y);}
 \draw[fill=orange!55, draw=orange!80!black] (0.8,0.8) rectangle (1.6,1.6);
 \node[anchor=north] at (1.2,-0.15) {局部卷积窗口};
 \draw[->, thick] (2.8,1.2)--(4.0,1.2);
 \draw[fill=green!10, draw=green!60!black] (4.4,0.2) rectangle (6.8,2.2);
 \draw[fill=green!45, draw=green!70!black] (5.1,0.9) rectangle (6.1,1.5);
 \node[anchor=north] at (5.6,0.05) {高层特征图};
 \node[above] at (5.6,2.35) {更大感受野、较低空间分辨率};
\end{tikzpicture}
\caption{卷积网络从局部窗口逐步形成层级表示。}
\label{fig:vision-receptive-field}
\end{figure}

图\ref{fig:vision-receptive-field}对照输入局部窗口与高层特征，突出扩大感受野同时降低空间分辨率的代价；细小缺陷是否被保留须在下采样前后检查。
\section{预训练、跨域泛化与工业验收}
标注稀缺时，监督微调并不是唯一选择。Contrastive Language--Image Pretraining (CLIP) 使用大规模图文对比学习把图像和文本映射到共享空间，给定“有裂纹”“无裂纹”或更细的缺陷描述即可进行零样本分类；但工业术语、相机域和细粒度缺陷通常与通用图文语料不同，因此 CLIP 只能作为迁移起点，必须通过领域文本模板、少量标注或 Low-Rank Adaptation (LoRA) 适配校准。Bootstrapping Language--Image Pretraining (BLIP)-2、LLaVA 等视觉语言模型把视觉编码器与语言模型连接起来，适合生成检验报告和解释性问答，但它们的文字流畅性不能替代像素级定位，报告中的缺陷位置仍应由检测器或分割器提供可验证证据。

\paragraph{自监督与视觉表征。}
当缺陷标签稀缺时，可以先用无标签图像学习表示。对比学习通过不同增强视图构造正样本，遮挡重建则要求模型恢复被隐藏的 patch。增强必须与工业语义一致：如果随机裁剪会删除缺陷，模型可能学会“正常背景”；如果两种视图来自不同设备，正样本关系也可能并不成立。预训练表示仍需在跨设备测试上验证，而不能只看线性探测分数。

\paragraph{校准与拒识。}
高置信度错误在质检系统中尤其危险。可以用温度缩放或保序回归修正分类概率，再根据风险设置拒识区域。拒识不是把困难样本删除，而是将其转交人工或更精细的模型。报告覆盖率--风险曲线可以说明系统在自动处理多少样本时保持怎样的错误水平。

\paragraph{任务边界与模型选择。}
分类回答“图像中是否存在某类状态”，检测还要回答“状态位于哪里”，分割则要求逐像素或逐区域描述。任务标签越细，标注成本越高，评价也越依赖边界和拓扑。若业务只需判断是否需要复检，像素级分割未必是最经济的方案；若需要估计裂纹长度和扩展方向，分类分数就不足以支持决策。

\paragraph{检测指标与匹配链。}
检测评价先按置信度排序预测框，再依据 IoU 和类别匹配真实框，形成精确率--召回率曲线。平均精度依赖匹配阈值、类别平均方式和置信度排序。小目标的一个像素偏移可能显著改变 IoU，因此应同时报告按目标尺寸分层的结果。非极大值抑制的阈值也会影响最终告警数量，不能只记录网络训练损失。

\paragraph{正负样本与困难样本。}
工业缺陷通常是少数类，检测器会面对大量容易分类的背景位置。焦点损失通过降低易样本权重，把梯度集中到困难样本；类别权重则直接改变不同类别的损失比例。两者都可能提高少数类召回，却也可能放大错误标注和难例噪声。验收时应同时查看每类召回、误报空间分布和人工复核负担。

\paragraph{分割后处理与几何指标。}
像素预测之后常需要连通域分析、孔洞填充、形态学操作或面积过滤。后处理可以改善视觉结果，却也可能删除细小真实缺陷。应把后处理规则纳入验证规则，并在不同缺陷尺寸、边界模糊程度和相机分辨率下检查稳定性。对裂纹等细长对象，骨架长度、断裂数和最大偏差可能比单一 Dice 更有解释力。

\paragraph{域偏移与测试时监控。}
光照、相机、镜头污染、材料批次和背景工装都会改变输入分布。输入统计监控可以发现均值、对比度和频谱变化，但不能直接证明标签关系改变。部署后应结合少量延迟标签检查分层性能，区分数据质量告警、需要校准的协变量偏移和需要重新标注的概念偏移。

\paragraph{解释工具的证据边界。}
梯度显著性、遮挡图和类激活图可以帮助定位模型关注的区域，但高亮处不一定就是导致故障判断的真实原因。例如模型可能看的是裂纹旁边的批次标记。可以分别遮住缺陷、替换背景或只保留缺陷区域，观察输出如何变化，并请专家核对。不过这些图像修改也可能产生不自然的输入，结果应结合跨相机测试解释，不能仅凭一张热力图作因果判断。

\paragraph{视觉系统验收。}
验收应覆盖正常样本、已知缺陷、少见缺陷、遮挡、模糊、光照变化、相机切换和未知输入。指标包括缺陷召回、每千张误报、定位偏差、处理延迟、拒识率和人工复核时间。若模型在实验室图像上表现优秀，却无法承受生产线速度和污染条件，就不能视为完成部署。

\paragraph{分层实验设置。}
设生产线上有四台相机、六个批次和三类缺陷。训练集可以使用前三台相机的早期批次，验证集使用相同相机的后期批次，测试集则完整留出第四台相机。这样的划分同时检验时间变化和设备迁移。若把所有图像随机打散，来自同一工件的近重复帧可能进入不同集合，检测结果会明显偏乐观。

实验应先建立三个层次的基线：多数类或人工规则基线、冻结预训练编码器的线性头、端到端微调模型。每次只改变一个因素，例如输入分辨率、增强策略、正负样本分配或后处理阈值。除 Mean Average Precision (mAP) 外，报告缺陷级召回、每千张正常图的误报、最小可检缺陷尺寸和单张图延迟。对生产线而言，漏掉一个关键缺陷的代价通常不能由平均准确率代表。

\paragraph{分割结果验收。}
对裂纹掩码，先计算 Dice 和 IoU，再计算边界到真实边界的平均距离、最大距离和骨架断裂数。若模型通过扩大预测区域提高召回，IoU 可能下降但安全风险降低；若模型只保留粗主干，Dice 可能尚可，却无法估计裂纹长度。因此验收规则应先由业务定义允许的漏检和过检，再选择指标和阈值。

标注不确定时，可以由多名标注者独立勾画并保留分歧区域。模型在分歧区域给出低置信度并不一定是失败，反而可能提示需要复核。训练、评价和部署中的掩码格式、像素坐标系、缩放方式必须保持一致，否则边界误差会被错误归因于模型。

\section{数据反馈与部署决策}
少样本适配不应只比较准确率。新相机只有少量标注时，可以冻结 ResNet、ConvNeXt、Swin 或 ViT 的大部分编码器，仅更新归一化统计量、分类头或低秩适配参数；也可以使用 CLIP 的文本原型进行零样本筛查，再把高不确定样本交给人工。验证时必须保留未参与适配的批次，并比较全量微调、参数高效微调和直接拒识三种方案。若新域置信度整体下降，拒识率应作为业务容量约束，而不是被隐藏的失败样本比例。

\paragraph{数据处理记录。}
每张图像应关联相机编号、镜头、曝光、生产批次、工件编号、标注版本和采集时间。增强后的图像保留原图引用，预测结果记录模型版本和预处理参数。发生误报时，只有根据这些记录回溯，才能判断问题来自光学条件、标注规则、数据切分、模型还是后处理。可追溯性是视觉智能进入生产系统的必要条件。

\paragraph{主动学习与标注预算。}
当人工标注昂贵时，可优先选择高不确定性、代表新域或具有高业务价值的图像。只按模型置信度采样可能反复选择相似噪声，因而需要结合特征多样性、设备覆盖和缺陷稀有度。每轮主动学习都应保留固定测试集，记录新增标注数量、类别分布、性能变化和人工时间，才能判断单位标注成本带来的真实收益。

以下材料按“残差与图像块表示—检测—像素分割”对应本章内容；卷积基础只保留一个复习入口。视频名称与链接已对照作者课程目录，未逐帧观看，播放量未能可靠核实，故不作热度排序。
\section*{辅助学习材料}
\begingroup
\emergencystretch=3em
\begin{itemize}
\item \textbf{Convolutional Neural Networks (CNNs / ConvNets)}（Stanford CS231n，英文笔记）。仅在需要复习时阅读卷积层、输出尺寸与参数共享，回到本章“卷积的索引、参数与等变性”检查步幅和填充。\href{https://cs231n.github.io/convolutional-networks/}{在线阅读}。
\item \textbf{ResNet网络讲解}（中文视频）。对照“残差梯度不再只有单一路径”，区分恒等捷径与改变尺寸的投影捷径，解释残差块为何比单纯加深网络更易优化。\href{https://www.bilibili.com/video/BV1T7411T7wa}{观看讲解}；中文博客\href{https://zhuanlan.zhihu.com/p/42706477}{《详解残差网络》}可作图文补充，其标题和残差块主题已由必应索引核对，但知乎原页访问受限，未核读全文。
\item \textbf{Vision Transformer网络讲解}（中文视频）。重点看图像如何切成 patch、投影为 token 并加入位置信息，再对照“图像块与窗口注意力的成本”计算分辨率变化带来的开销；这是 ViT 的图像建模讲解，不是泛自注意力入门。\href{https://www.bilibili.com/video/BV1Jh411Y7WQ}{观看讲解}。
\item \textbf{区域卷积神经网络（R-CNN）系列}（《动手学深度学习》，中文教程）。先读其中的 Fast R-CNN 小节与 RoI pooling 示例，将共享整图特征、候选区域、分类和框回归对应到“检测框、匹配与定位损失”。Fast R-CNN 使用选择性搜索；随后对照 Faster R-CNN 小节，辨明后者才引入区域提议网络（RPN），不可混称。\href{https://zh.d2l.ai/chapter\_computer-vision/rcnn.html}{在线阅读}。
\item \textbf{YOLO系列网络讲解(V1\textasciitilde V3)}（中文视频）。沿单阶段预测到多尺度检测的变化，比较网格、锚框与类别／位置输出，对照本章“空间输出”和“检测匹配与正样本分配”；不要把早期版本的机制直接套给所有新版 YOLO。\href{https://www.bilibili.com/video/BV1yi4y1g7ro}{观看讲解}。
\item \textbf{U-Net网络讲解}（中文视频）。沿编码器、解码器与同尺度特征拼接追踪分辨率，回到“语义分割、实例分割与 U-Net”检查细裂纹边界为何需要低层特征；并与 ResNet 的相加捷径区分。\href{https://www.bilibili.com/video/BV1Vq4y127fB/}{观看讲解}。上述四个模型视频均可从\href{https://github.com/WZMIAOMIAO/deep-learning-for-image-processing}{作者课程目录}查到配套代码。
\item \textbf{TorchVision Object Detection Finetuning Tutorial}（PyTorch 官方，英文实践）。跟随自定义数据集、预训练 Mask R-CNN 微调与评价流程，核对框、标签和掩码格式，对应本章“预训练、跨域泛化与工业验收”；它是检测／实例分割实践，不替代前面的 Fast R-CNN 原理解说。\href{https://docs.pytorch.org/tutorials/intermediate/torchvision\_tutorial.html}{在线阅读}。
\end{itemize}
\endgroup
\paragraph{本章小结。}
视觉表示的选择应从输出要求出发：分类需要区分状态，检测需要定位区域，分割还要恢复边界。卷积和残差提供层级特征，注意力扩展区域之间的信息交互，预训练则改变这些表示所需的标注预算。最终方案必须同时满足缺陷尺度、跨域可靠性、推理资源和人工复核容量的约束。卷积索引与感受野决定证据是否仍在特征中，残差路径决定这些特征能否有效训练，框回归、焦点损失和Dice则分别改变位置误差、难例与区域重叠的权重。任何一个损失下降都只回答其中一部分问题，不能代替空间匹配、跨相机验证和部署后的复核。

找到裂纹后，还得查这台设备允许多大缺陷、该按哪一版规程复检。图像本身不能回答这些问题，下一章因此讨论语言模型怎样检索和使用手册、工单。生成报告时应保留裂纹坐标、尺寸尺度、拍摄时间和模型版本，不能只用一句“发现异常”替代可复查的检测结果。
